\clearpage
\thispagestyle{plain}
\section*{Worum es in diesem Buch geht}

Dieses Buch geht auf das Werk \emph{ad-Durūs al-Muhimmah li-ʿĀmmati l-Ummah} von Shaykh
ʿAbd al-ʿAzīz ibn ʿAbdillāh ibn Bāz \ar{رحمه الله} zurück. Der Shaykh selbst beschrieb sein Werk
als knappe Worte zur Erläuterung einiger Informationen, die jeder Muslim über den Islam
erfahren soll.

Warum beschäftigen wir uns gerade mit diesen Lektionen? Weil sie wichtig sind – so wie es der
Verfasser selbst bezeichnete und wie es die Gelehrten empfahlen. Und sollte jemand meinen, er
stehe über diesen Grundlagen, weil er bereits weiter fortgeschritten sei, dann gilt: Man frage
ihn nach dem Inhalt dieser Lektionen. Kennt er ihn nicht sicher, ist die Allgemeinheit der
Muslime, die ihn kennt, besser dran als er. Der Student soll bescheiden bleiben, sich gegenüber
dem Wissen und den Gelehrten nicht überheblich verhalten und dem Weg der rechtschaffenen
Gelehrten folgen.

\textbf{Dieses Buch behandelt:}
\begin{enumerate}[leftmargin=1.6em]
  \item den Umgang der rechtschaffenen Vorfahren mit dem Qurʾān – wie sie ihn rezitierten,
        auswendig lernten, studierten und umsetzten,
  \item die Grundsätze von Hingabe (\emph{islām}), Glauben (\emph{īmān}) und Pietät
        (\emph{iḥsān}), den Glauben an die Einzigkeit Allahs (\emph{tawḥīd}) und die
        verschiedenen Arten der Beigesellung (\emph{širk}),
  \item das Gebet (\emph{ṣalāh}),
  \item die rituelle Gebetswaschung (\emph{wuḍūʾ}),
  \item die Verpflichtung zu rechtschaffenen Sitten und Moralvorstellungen,
  \item die Warnung vor Beigesellung und verschiedenen Sünden,
  \item die rituelle Bestattung des Verstorbenen, das Totengebet und das Begräbnis.
\end{enumerate}

Zur inhaltlichen Erläuterung stütze ich mich in diesem Buch durchgehend auf den Šarḥ (die
ausführliche Erklärung) von Shaykh Hayṯam Ibn Muḥammad Ğamīl Sarḥān, ehemaliger Dozent
am Ḥaram-Institut an der Propheten-Moschee in Medina. Ich lege hier keinen eigenen,
selbstständigen Šarḥ vor, sondern orientiere mich in Aufbau und Inhalt eng an seiner
Erläuterung und bereite sie für den gemeinsamen Unterricht auf.

Möge Allah \ar{ﷻ} uns allen nützliches Wissen, Aufrichtigkeit und Beständigkeit schenken.
\clearpage
